% Sección: Parte 1 – Problemas físicos complejos

% Frame principal de la sección
\begin{frame}{Parte 1: Problemas físicos complejos}
  \begin{itemize}
    \item Muchos fenómenos en física están descritos por ecuaciones diferenciales parciales no lineales de alta dimensión.
    \item Algunos ejemplos clásicos:
      \begin{itemize}
        \item Dinámica de fluidos (turbulencia, flujo transitorio)
        \item Ecuaciones de Navier–Stokes en geometrías complejas
        \item Problemas de transporte radiativo en plasmas
        \item Mecánica cuántica de muchos cuerpos (Hamiltonianos de alta dimensión)
        \item Sistemas caóticos y sensibles a condiciones iniciales
      \end{itemize}
    
  \end{itemize}
\end{frame}


\begin{frame}{Parte 1: Problemas físicos complejos}
	\begin{itemize}
		\item Retos comunes:
      \begin{itemize}
        \item Altísimo costo computacional en mallado fino
        \item Inestabilidad numérica y rigidez del problema
        \item Necesidad de grandes recursos (CPU/GPU) y tiempo de cómputo
        \item Dificultad para generalizar soluciones: recalcular desde cero ante cambios de parámetro
      \end{itemize}
    \item Conclusión: para estos problemas, los métodos tradicionales (dif. finitas, elementos finitos) pueden quedarse cortos. \alert{Se requiere una aproximación novedosa.}
	\end{itemize}
\end{frame}
